%% COULD BE USEFUL LATER ON
%\usepackage[a4paper,top=3cm,bottom=2cm,left=3cm,right=3cm,marginparwidth=1.75cm]{geometry}
%\usepackage[a-1b]{pdfx}   % for PDF/A-1b - Submit thesis

%%%NON SO SE E' LA GEOMETRY CHE MI SERVE!!!!
%\newgeometry{
	%top=3cm,
	%bottom=2.5cm,
	%outer=2cm,
	%inner=2.9cm,
	%}
	
%%%%%%%%%%%%%%%%%%%%%%%%%%%%%%%%%%%%%%%%%%%%%%%%%%%%%%%%%%%%%%%%%%%%%%%%%%%%%%%%%%%%%%%%%%%%%%%%%%%%%%%%%%%%%%%%%%%%%%%%
%%%%%%%%%%%%%%%%%%%%%%%%%%%%%%%%%%%%%%%%%%%%%%%%%%%%%%%%%%%%%%%%%%%%%%%%%%%%%%%%%%%%%%%%%%%%%%%%%%%%%%%%%%%%%%%%%%%%%%%%
	

%% Useful packages
\usepackage{dsfont}
\usepackage{amsmath}
\usepackage{amsthm}
\usepackage{amsfonts}
\usepackage{amssymb}
\usepackage{braket} % se usi \ket, \bra...
\usepackage{bbm}\usepackage{bbold}
\usepackage{mathtools} %Symbol of definition
\usepackage{graphicx} % Required for inserting images
\usepackage{wrapfig} % Per poter piazzare figure immerse nel testo
\usepackage{marvosym} % Per aggiungere il punto d'ariete come \Aries
\usepackage{verbatim} % Multiline comments con \begin{comment} ... \end{comment}
\usepackage{comment}
\usepackage{url}
\usepackage[colorinlistoftodos]{todonotes}
\usepackage[colorlinks=true, allcolors=blue]{hyperref}
\usepackage{chngcntr}
\usepackage{breqn} % Spezzare le equazioni
\usepackage{siunitx} %Units of measure
\usepackage{physics}  % \ket, \bra, \outerproduct, \comm, \dv, ...
\usepackage{tikz}
\usetikzlibrary{arrows.meta,positioning,decorations.pathreplacing}

%% Structure

\newenvironment{abstract}{}{}  % definizione minima, vuota
\usepackage{abstract}

\usepackage{caption}
\usepackage{subcaption}
\usepackage{frontespizio}
\usepackage{multicol}
\usepackage{layouts}
\usepackage{tocloft}
\usepackage{titlesec}
\usepackage{etoolbox,refcount}

\titleformat
{\chapter} % command
[display] % shape
{\bfseries\Large} % format
{\thechapter} % label
{0.5ex}
{}
%c++
\usepackage{listings}
\usepackage{xcolor}
\usepackage{minted}

% Images and bibliography
\graphicspath{ {./Images/} }
\usepackage{filecontents}
\usepackage[backend=biber, style=numeric,sorting=none]{biblatex}
\addbibresource{Bibliography.bib}


%%%%%%%%%%%%%%%%%%%%%%%%%%%%%%%%%%%%%%%%%%%%%%%%%%%%%%%%%%%%%%%%%%%%%%%%%%%%%%%%%%%%%%%%%%%%%%%%%%%%%%%%%%%%%%%%%%%%%%%%
%%%%%%%%%%%%%%%%%%%%%%%%%%%%%%%%%%%%%%%%%%%%%%%%%%%%%%%%%%%%%%%%%%%%%%%%%%%%%%%%%%%%%%%%%%%%%%%%%%%%%%%%%%%%%%%%%%%%%%%%

\usepackage[english]{babel}
\usepackage[utf8]{inputenc}
\usepackage[document]{ragged2e}
\usepackage[T1]{fontenc}


%%%%%%%%%%%%%%%%%%%%%%%%%%%%%%%%%%%%%%%%%%%%%%%%%%%%%%%%%%%%%%%%%%%%%%%%%%%%%%%%%%%%%%%%%%%%%%%%%%%%%%%%%%%%%%%%%%%%%%%%
%%%%%%%%%%%%%%%%%%%%%%%%%%%%%%%%%%%%%%%%%%%%%%%%%%%%%%%%%%%%%%%%%%%%%%%%%%%%%%%%%%%%%%%%%%%%%%%%%%%%%%%%%%%%%%%%%%%%%%%%

%%Inserisco il file 
% =========================
% Frequently-used shortcuts
% =========================



%%%%%%%%%%%%%%%%%%%%%%%%%%%%%%%%%%%%%%%%%%%%%%%%%%%%%%%%%%%%
%Old-Commands
\newcommand{\mbeq}{\overset{!}{=}}
\newcommand\lreqn[2]{\noindent\makebox[\textwidth]{$\displaystyle#1$\hfill(#2)}\vspace{2ex}}
\newcommand*{\bigchi}{\mbox{\Large$\chi$}}% big chi
\newcommand\numberthis{\addtocounter{equation}{1}\tag{\theequation}}
\newcommand{\D}[1]{\mathcal{D}\!\left[#1\right]}
\newcommand{\liouvillian}{\mathcal{L}}  %Liouvilliano
%%%%%%%%%%%%%%%%%%%%%%%%%%%%%%%%%%%%%%%%%%%%%%%%%%%%%%%%%%%%

%Collapsing operators
\newcommand{\Cp}{\hat C_+}
\newcommand{\Cm}{\hat C_-}

\newcommand{\Cpdag}{\hat C_+^\dagger}
\newcommand{\Cmdag}{\hat C_-^\dagger}

\newcommand{\Cpdnd}{\hat C_+^\dagger\hat C_+}
\newcommand{\Cmdnd}{\hat C_-^\dagger\hat C_-}

% Identity and time step
\newcommand{\id}{\mathds{1}}
\newcommand{\dt}{\delta t}
\newcommand{\Iddt}{\id\,\dt}

% Pauli (generic)
\newcommand{\sx}{\sigma_x}
\newcommand{\sy}{\sigma_y}
\newcommand{\sz}{\sigma_z}

% Pauli on qubit i
\newcommand{\sxi}[1]{\sigma_x^{(#1)}}
\newcommand{\syi}[1]{\sigma_y^{(#1)}}
\newcommand{\szi}[1]{\sigma_z^{(#1)}}

% Collective spin operators (hatted)
\newcommand{\Jx}{\hat J_x}
\newcommand{\Jy}{\hat J_y}
\newcommand{\Jz}{\hat J_z}
\newcommand{\Jp}{\hat J_{+}}
\newcommand{\Jm}{\hat J_{-}}

% Two-qubit parity + projectors
\newcommand{\Piz}{\Pi_z}
\newcommand{\Pizdef}{\szi{1}\szi{2}}
\newcommand{\Pp}{\Pi_{+}}
\newcommand{\Pm}{\Pi_{-}}
\newcommand{\Ppdef}{\frac{1}{2}\qty(\id+\Piz)}
\newcommand{\Pmdef}{\frac{1}{2}\qty(\id-\Piz)}

% Effective Hamiltonian and superoperators
\newcommand{\Heff}{\hat H_{\mathrm{eff}}}
\newcommand{\Liou}{\mathcal{L}}
\newcommand{\Diss}[1]{\mathcal{D}\!\qty[#1]}      % dissipator as a map
\newcommand{\Dact}[2]{\mathcal{D}\!\qty[#1]\,#2} % action on an operator/state

% Tilde ket that does NOT stretch the delimiters
\newcommand{\tket}[1]{\ket{\smash[t]{\tilde{#1}}}}

% CSS shorthand (your notation)
\newcommand{\CSS}[2]{\ket{#1,#2}}
\newcommand{\CSSprod}[3]{%
	\bigotimes_{i=1}^{#3}\left(
	e^{i #2}\sin\frac{#1}{2}\ket{1}_i+\cos\frac{#1}{2}\ket{0}_i
	\right)}

% Spin-squeezing symbols (handy labels)
\newcommand{\xiKU}{\xi_{\mathrm{KU}}^2}
\newcommand{\xiW}{\xi_{\mathrm{W}}^2}



% Header/footer
\usepackage{fancyhdr}



\fancypagestyle{plain}{%
	\fancyhf{}%
	\renewcommand{\headrulewidth}{0pt}%
	\fancyhead[LE,RO]{\thepage}%
}


\fancypagestyle{cleared}{%
	\fancyhf{}            % svuota header/footer
	\fancyfoot[C]{\thepage} % numero pagina centrato in basso
	\renewcommand{\headrulewidth}{0pt} % niente linea in alto
	\renewcommand{\footrulewidth}{0pt} % niente linea in basso
}



% Se usi biblatex:
% \usepackage[backend=biber,style=phys,sorting=none]{biblatex}
% \addbibresource{bibliography.bib}
